\section{Control, configuration and monitoring protocols}

\subsection{Service architecture}

SC provides the OPC UA server and control arbitration. The server hosts DAQ
configuration and forwards publications from
\code{daphne-server}: equipment monitoring to Ignition; firmware counters,
timing state and applied readout configuration to DAQ opmon.
Hermes receives configuration over UDP/IPBus. Figure~\ref{fig:protocol} shows
the DAPHNE path; box interfaces are defined in Section~\ref{sec:boxes}.

\begin{figure}[htbp]
\centering
\begin{tikzpicture}
\node[box,text width=29mm] (ignition) at (0,0) {Ignition\\SC client};
\node[box,text width=32mm] (daq) at (5.0,0) {DAQ CCM client\\opmon subscriber};
\node[box,text width=35mm] (hermesclient) at (10.7,0)
 {DAQ Hermes\\configuration client};
\node[box,text width=75mm,minimum height=19mm] (ua) at (2.5,-2.3)
 {External OPC UA server\\Integrated DAQ configuration\\Protobuf / ZeroMQ client and subscriber};
\node[box,text width=49mm,minimum height=18mm] (daphne) at (2.5,-5.8)
 {Board-side daphne-server\\Configuration endpoint\\Monitoring publisher};
\node[box,text width=37mm,minimum height=18mm] (hermes) at (10.7,-5.8)
 {Board-side Hermes\\IPBus / UDP\\Configuration endpoint};
\node[draw=dunegreen,dashed,rounded corners,fit=(daphne)(hermes),inner sep=8pt,
 label={[font=\footnotesize,text=dunegreen]below:One DAPHNE board; one shared 1\,Gb Ethernet CCM connection}] {};
\draw[both] (ignition) -- node[left,font=\footnotesize]{OPC UA} (ignition |- ua.north);
\draw[both] (daq) -- node[right,font=\footnotesize]{OPC UA} (daq |- ua.north);
\draw[both] ([xshift=-11mm]ua.south) -- node[left,align=right,font=\footnotesize]
 {Configuration\\request/response} ([xshift=-11mm]daphne.north);
\draw[flow,draw=dunegreen] ([xshift=11mm]daphne.north) -- node[right,align=left,font=\footnotesize,fill=white,inner sep=2pt]
 {Published\\monitoring} ([xshift=11mm]ua.south);
\draw[both] (hermesclient.south) -- node[right,align=left,font=\footnotesize]
 {IPBus / UDP\\configuration} (hermes.north);
\end{tikzpicture}
\caption{DAPHNE CCM configuration and published monitoring.}
\label{fig:protocol}
\end{figure}

\subsection{Configuration exchanges}

DAQ CCM requests run-variable and timing-endpoint configuration through OPC UA.
The server exchanges Protobuf
messages over ZeroMQ and TCP using a DEALER client and the board's ROUTER
endpoint [R3, R4].
The reference port is TCP 40001 [R5]. Correlated responses shall distinguish
acceptance, completion and failure.

Hermes retains its independent IPBus/UDP configuration endpoint, reference
UDP 50001 [R6], outside the OPC UA server. Both services share the board
connection and gateware lifecycle.

\subsection{Published monitoring}

Monitoring shall be published periodically or on change, without a client
request for each update. Only configuration and control operations, including
reset/recovery commands, use application request/response. Local hardware
sampling is performed by the producing service.

The proposed DAPHNE publication transport is Protobuf over ZeroMQ PUB/SUB over TCP,
using endpoints separate from configuration. The OPC UA server subscribes to
\code{daphne-server} publications and distributes them through OPC UA
subscriptions. Opmon monitoring follows the path
\code{daphne-server} $\rightarrow$ OPC UA server $\rightarrow$ DAQ opmon.

\begin{center}\small
\begin{tabular}{L{0.25\textwidth}L{0.68\textwidth}}
\toprule
\textbf{Consumer} & \textbf{Monitoring responsibility} \\
\midrule
DAQ opmon & Firmware counters, timing state and applied data-readout configuration, published by \code{daphne-server} through OPC UA. \\
SC / Ignition & Voltages, temperatures, fans, service status and software/firmware versions. \\
SC and DAQ & Timing-endpoint condition, reset/recovery progress and board readiness. \\
\bottomrule
\end{tabular}
\end{center}

Opmon shall not acquire voltage monitoring. SC receives equipment monitoring
through the OPC UA server without an opmon dependency. Publications shall
carry source identity, time, sequence and validity. Subscribers shall detect
stale or missing updates and obtain a current snapshot after reconnecting.

\subsection{Timing-endpoint control and recovery over GbE}

Timing-endpoint configuration follows the DAQ configuration path:
\begin{center}
DAQ $\longleftrightarrow$ OPC UA server $\longleftrightarrow$
DAPHNE service $\longleftrightarrow$ local timing endpoint.
\end{center}
The same interface shall expose reset and recovery to authorized SC clients.
Timing monitoring shall be published over the board's GbE connection and
made available to both SC and DAQ.

A reset request shall identify its target and receive an execution response.
The board shall publish reset assertion, release and the resulting endpoint
state so that both clients can acknowledge the same recovery operation.
Recovery access shall remain available while timing is unavailable. Concurrent
DAQ and SC commands shall be serialized by SC's control infrastructure.

SC shall determine and publish \emph{board ready for data taking} from the
agreed readiness conditions, including completed reset, restored clocks,
endpoint readiness and valid synchronization [R4]. Reset release alone is
insufficient. Missing or stale required publications shall withhold readiness.
DAQ consumes the published readiness before resuming data taking.

\subsection{PDS Calibration Box and PoF laser boxes}
\label{sec:boxes}

Each box controller shall expose configuration/control request/response and
published monitoring through its native service and an adapter in the external
OPC UA server. Ignition subscribes to equipment monitoring; DAQ subscribes to
the readiness and operating status needed for data taking.

\textbf{Calibration Box (LCM).} PDS supplies the calibration hardware, firmware
and service. DAQ uses that service to request calibration settings and sequence
execution. SC provides command arbitration, equipment power/service control
and monitoring of health, alarms and calibration status. The OPC UA interface
shall expose protective inhibit actions.

\textbf{PoF laser boxes.} SC controls optical-power operation and receives
controller health, operating state, permits and inhibit status. The interface
shall expose DPS protection status and give protective inhibits priority over
operating commands. PDS and DPS shall define the controller/DPS connection and its
acknowledgement behavior.

PDS shall publish each controller's native protocol, transport, ports, command
acknowledgements, monitoring schema and compatible versions. These details
remain to be specified for both box types; the DAPHNE Protobuf/ZeroMQ contract
applies only to DAPHNE.

\clearpage
\subsection{Interface release and verification}

Each release shall publish compatible versions and interface specifications
online: schemas, OPC UA model, Hermes and box-controller interfaces, and
endpoint/port allocation. Contracts shall define authentication, timeouts,
queue limits, update cadence and reconnect behavior. DAPHNE publication
schemas and ports remain to be assigned.

Verification shall cover concurrent configuration and publication on the shared
link, the HD and VD installation scales, subscriber loss/reconnection, and SC
timing recovery with a board-generated completion acknowledgement. Box
verification shall cover command ownership, published status, reconnect
behavior and inhibit acknowledgement. The
DAQ--PDS and DAQ--SC/DPS interface documents shall use this same CCM contract.
